\documentclass[conference]{IEEEtran}

\usepackage{amsmath,amssymb}
\usepackage{booktabs}
\usepackage{balance}
\usepackage{cite}
\usepackage{graphicx}
\usepackage{microtype}
\usepackage[hidelinks]{hyperref}

\newcommand{\method}{SC-ICDM}
\newcommand{\vect}[1]{\boldsymbol{#1}}
\newcommand{\norm}[1]{\left\lVert#1\right\rVert}
\newcommand{\pos}[1]{\left[#1\right]_{+}}
\hyphenation{com-mu-ni-ca-tions}

\begin{document}

\title{Self-Calibrated Diffusion Interference Cancellation\\
under Nonstationary Interference}

\author{
\IEEEauthorblockN{Author Name(s)}
\IEEEauthorblockA{
Department or Institution\\
City, Country\\
email@domain.com}
}

\maketitle

\begin{abstract}
Interference cancellation diffusion models (ICDMs) can recover semantic
features by combining separate diffusion priors for the desired signal and
the interference. Existing ICDM inference, however, requires the receiver to
know a single frame-level signal-to-interference-plus-noise ratio (SINR).
That assumption is restrictive when the interference power is unknown and
changes within a frame. We develop a SINR-blind receiver that estimates
interference power from the received latent sequence and refines blockwise
amplitudes during reverse diffusion. The refinement reuses the current
denoised signal and interference estimates and therefore requires no
additional neural-network evaluations. In experiments with stationary,
alternating, and burst interference, blind global calibration remains within
0.008~dB of the true-global-SINR reference, while one-shot blockwise
calibration gains 0.819~dB and 2.108~dB in the two nonstationary settings.
Alternating calibration adds 0.599~dB over the one-shot blockwise receiver
(95\% image-pair-cluster bootstrap interval: 0.551--0.648~dB). The gain is
not universal: alternating updates lose 0.109--0.174~dB under the strongest
burst interference and remain 1.070~dB below direct decoding when
interference is weak. These results show that reconstruction-optimal
self-calibration need not recover the physical interference power and
motivate reliability-gated activation rather than unconditional diffusion
cancellation.
\end{abstract}

\begin{IEEEkeywords}
semantic communications, diffusion models, interference cancellation,
blind calibration, nonstationary interference
\end{IEEEkeywords}

\section{Introduction}

Diffusion priors provide a natural way to exploit the structure of
transmitted semantic features at the receiver. In the recently proposed
interference cancellation diffusion model (ICDM)~\cite{wu2026icdm}, one
diffusion model represents the desired latent distribution and a second
model represents the interference distribution. A shared likelihood term
couples the two reverse processes, allowing the receiver to reconstruct the
desired feature while explaining structured interference. This design
substantially improves reconstruction when the interference statistics used
by the receiver match the channel.

One piece of side information remains easy to overlook: the published
receiver selects the interference amplitude and guidance coefficients from
the true frame-level SINR. In deployment, this quantity may be unavailable.
More importantly, a single scalar cannot describe interference whose power
changes within a frame. Replacing the true SINR with a global energy estimate
addresses the first issue but not the second. Estimating each block
independently captures local variation, but short-block estimates are noisy
and are computed before the diffusion priors have separated the desired and
interfering components.

This paper asks whether an ICDM receiver can calibrate itself from the
received latent sequence, without access to the true SINR, when interference
power varies within a frame. We answer this question with a receiver-side
procedure that first obtains blind blockwise power estimates and then updates
the associated amplitudes from the denoised signal and interference estimates
produced during reverse diffusion. The update closes a feedback loop between
reconstruction and calibration without retraining either prior or adding
neural-network evaluations.

The main contributions are as follows:
\begin{itemize}
    \item We formulate a SINR-blind ICDM protocol in which the receiver
    observes only the noisy latent sequence, noise level, desired-link channel
    state, and a fixed block partition. Simulated SINR, clean latents, and
    true block powers are excluded from the proposed receiver.
    \item We introduce blockwise alternating calibration inside reverse
    diffusion. The current denoised components update local interference
    amplitudes and the corresponding guidance coefficients.
    \item We isolate the value and limits of self-calibration through matched
    comparisons. Blind global calibration preserves the mean performance of
    the true-global-SINR reference; local calibration improves nonstationary
    cases; and alternating feedback adds 0.599~dB over one-shot blockwise
    calibration. We also identify severe-burst and weak-interference regimes
    in which unconditional updates or cancellation are undesirable.
\end{itemize}

The scope is deliberately limited. The experiments use additive white
Gaussian noise (AWGN), a known block partition, a known noise level, and
fixed desired and interference priors. The receiver is therefore SINR-blind,
not fully blind. The evidence does not establish robustness to fading,
unknown interference families, or unknown change points.

\section{Related Work}

Diffusion models have been used as learned priors for wireless denoising and
semantic reconstruction. CDDM applies a diffusion prior after channel
equalization to suppress channel noise~\cite{wu2024cddm}, whereas ICDM
introduces separate signal and interference priors and couples their reverse
processes through the observation model~\cite{wu2026icdm}. We retain the
frozen ICDM priors and sampler but remove the true-SINR input and allow
interference power to vary across blocks.

Blind diffusion inverse problems also estimate unknown forward operators or
nuisance parameters. GibbsDDRM alternates posterior sampling with inference
of an unknown linear operator~\cite{murata2023gibbsddrm}. Moment-projected
diffusion methods incorporate approximate conditional covariance information
into inverse-problem guidance~\cite{boys2024tweedie}. These works establish
that parameter uncertainty and higher-order information are not new in
diffusion inference. Our contribution is a lightweight point-calibration
mechanism and a controlled receiver protocol for structured, intra-frame
interference in semantic communication; we make no posterior-calibration
claim.

\section{System Model}

\subsection{Blockwise Interference}

An image is encoded into $N$ complex symbols
$\vect{x}=f_{\mathrm{enc}}(\vect{s})$, normalized such that
$\norm{\vect{x}}_2^2/N=1$. An independently generated interference image is
encoded into $\vect{z}$ with the same framewise normalization. A fixed
partition divides the transmitted order into disjoint blocks
$\{\mathcal{B}_b\}_{b=1}^{B}$. For AWGN, the observation in block $b$ is
\begin{equation}
    \vect{y}_b=\vect{x}_b+a_b\vect{z}_b+\vect{n}_b,\qquad
    \vect{n}_b\sim\mathcal{CN}(\vect{0},N_0\vect{I}),
    \label{eq:system}
\end{equation}
where $a_b\geq 0$ is an unknown interference amplitude. The nominal
frame-level interference power is
\begin{equation}
    p_{\mathrm{int}}=
    \frac{1}{N}\sum_{b=1}^{B}|\mathcal{B}_b|a_b^2.
    \label{eq:nominal_power}
\end{equation}
The simulated SINR is
$10\log_{10}[1/(p_{\mathrm{int}}+N_0)]$. Because both latents are normalized
over the complete frame rather than within each block, the realized block
energy need not equal $a_b^2$.

We evaluate three profiles with identical nominal frame power. The stationary
profile uses the same amplitude in every block. The alternating profile
repeats normalized power multipliers $(0.25,1.75)$, and the burst profile
repeats $(0,2)$. Real and imaginary components of one complex symbol share
the same block coefficient.

\subsection{Receiver Information}

The receiver knows $\vect{y}$, $N_0$, the block partition, and the two frozen
diffusion priors. It does not observe $\vect{x}$, $\vect{z}$, $a_b$, the
power-profile label, or the simulated SINR. The block partition is assumed
known; detecting change points is outside the present scope. True SINR is
used only to generate and stratify test examples.

\section{Self-Calibrated ICDM}

\subsection{Blind Initialization}

Under the unit-power and independence assumptions, the interference power in
block $b$ is initialized by subtracting the expected signal and noise powers
from the received energy:
\begin{equation}
\begin{split}
    \widehat{p}_b^{(0)}
    &=
    \pos{
    \frac{1}{|\mathcal{B}_b|}
    \sum_{k\in\mathcal{B}_b}|y_k|^2-1-N_0},\\
    \widehat{a}_b^{(0)}
    &=\sqrt{\widehat{p}_b^{(0)}}.
\end{split}
\label{eq:initialization}
\end{equation}
The blind global receiver applies~\eqref{eq:initialization} to the complete
frame; the blind blockwise receiver applies it separately to each block.
Each estimated power gives a local SINR
\begin{equation}
    \widehat{\gamma}_b
    =-10\log_{10}\!\left(\widehat{p}_b+N_0\right),
    \label{eq:estimated_sinr}
\end{equation}
which selects the ICDM likelihood-guidance coefficients by linear
interpolation of the original lookup table. Values outside the table are
clipped to its endpoints. Neither~\eqref{eq:initialization}
nor~\eqref{eq:estimated_sinr} accesses the true SINR.

\subsection{Diffusion-in-the-Loop Update}

One-shot energy estimation ignores information revealed as the two diffusion
priors separate the received mixture. At reverse step $t$, the pretrained
noise predictors give clean-point estimates
\begin{equation}
\begin{split}
    \widehat{\vect{x}}_0^{(t)}
    &=
    \frac{\vect{x}_t-\sigma_t
    \vect{\epsilon}_{\theta_x}(\vect{x}_t,t)}{\alpha_t},\\
    \widehat{\vect{z}}_0^{(t)}
    &=
    \frac{\vect{z}_t-\sigma_t
    \vect{\epsilon}_{\theta_z}(\vect{z}_t,t)}{\alpha_t}.
\end{split}
\label{eq:denoised}
\end{equation}
We project each complete-frame estimate to unit complex power, consistent with
the transmitter normalization, and denote the projected estimates by
$\widetilde{\vect{x}}_0^{(t)}$ and $\widetilde{\vect{z}}_0^{(t)}$. Once
$\alpha_t\geq0.5$, \method{} updates each block by nonnegative least squares:
\begin{equation}
    \widehat{a}_b^{(t)}
    =
    \pos{
    \frac{
    \operatorname{Re}\!
    \left\langle
    \widetilde{\vect{z}}_{0,b}^{(t)},
    \vect{y}_b-\widetilde{\vect{x}}_{0,b}^{(t)}
    \right\rangle}
    {\max\!\left(
    \norm{\widetilde{\vect{z}}_{0,b}^{(t)}}_2^2,\tau_b
    \right)}}.
    \label{eq:update}
\end{equation}
Here $\tau_b=|\mathcal{B}_b|\epsilon_{\mathrm{machine}}$. If the
denominator does not exceed $\tau_b$, the preceding amplitude is retained.
The updated power changes both the interference scale and the interpolated
guidance coefficients at the next correction. With 40 order-3 UniPC steps,
the schedule performs 21 updates.

Equation~\eqref{eq:update} reuses model outputs that ICDM already computes, so
it introduces no additional denoiser call. It is nevertheless a heuristic
point update: the framewise projection is not a posterior-moment identity,
and $\widehat{a}_b^{(t)}$ should not be interpreted as a calibrated
measurement of physical interference power.

\subsection{Compared Receivers}

The experiments compare direct JSCC decoding without diffusion cancellation;
the global oracle (A0), which receives the true frame-level SINR; blind global
calibration (A2); one-shot blind blockwise calibration (A3); and the proposed
alternating receiver \method{} (A4). We also evaluate the power-convention
correction A1 as an implementation control. A0 is a global-parameter
reference, not a performance upper bound: it knows the correct frame average
but not the nonstationary block powers.

\section{Experimental Setup}

We use 200 CelebA test images as desired sources and pair them with CIFAR-10
interference images. The SwinJSCC encoder and decoder, the interference
encoder, and both DiT priors are fixed for every receiver. Desired images are
reconstructed at $128\times128$ resolution. The AWGN SNR is 20~dB, and the
simulated global SINR belongs to $\{-7,-4,0,4,7,10\}$~dB. Each image pair is
evaluated with three channel seeds, all three power profiles, and all six
operating points. The latent block length is 64 complex symbols. Every
diffusion receiver uses 40 order-3 UniPC steps and identical model
evaluations.

The baseline run contains 54,000 reconstructions across Direct and A0--A3;
the matched \method{} run adds 10,800 reconstructions. For each comparison,
the image, interference image, channel realization, and diffusion
initialization are identical. Checkpoint and source hashes are retained with
the raw outputs.

We report PSNR and paired differences. Confidence intervals for \method{}
contrasts use 10,000 percentile bootstrap resamples of the 200 image-pair
clusters, keeping every seed, SINR, and profile belonging to an image pair in
the same resample. The image pair is the experimental unit; treating the
10,800 repeated records as independent would understate uncertainty. The
current 200 pairs were inspected during method development, so the intervals
in this draft are exploratory and not multiplicity-adjusted.

Receiver time includes calibration, diffusion sampling, output normalization,
and JSCC decoding, with GPU synchronization. Data loading, source encoding,
and channel simulation are excluded.

\section{Results}

\subsection{Global and Local Blind Calibration}

The blind global receiver tracks the global oracle closely. Its mean
difference from A0 is $-0.005$~dB under stationary interference, $-0.008$~dB
under alternating interference, and $-0.005$~dB under burst interference.
These descriptive gaps do not prove statistical equivalence, but they show
that true receiver-side SINR is unnecessary for reproducing the mean
performance of the published global receiver in this setup. Correcting the
original amplitude convention changes mean PSNR by only
0.007--0.010~dB and therefore does not explain the gains below.

Local estimation matters when interference is nonstationary. Relative to A0,
the one-shot blockwise receiver loses 0.292~dB in the stationary case but
gains 0.819~dB under alternating interference and 2.108~dB under burst
interference. A correct global average is therefore insufficient when power
variation within the frame carries useful receiver information.

\begin{table}[t]
\caption{Mean Reconstruction PSNR (dB)}
\label{tab:mean_psnr}
\centering
\setlength{\tabcolsep}{3.3pt}
\begin{tabular}{lrrrrr}
\toprule
Profile & Direct & A0 & A2 & A3 & \method{} \\
\midrule
Stationary  & 18.616 & 19.058 & 19.053 & 18.767 & \textbf{19.446} \\
Alternating & 18.811 & 18.939 & 18.931 & 19.758 & \textbf{20.352} \\
Burst       & 19.047 & 18.903 & 18.898 & 21.012 & \textbf{21.535} \\
\midrule
All         & 18.825 & 18.967 & 18.961 & 19.845 & \textbf{20.444} \\
\bottomrule
\end{tabular}
\end{table}

\subsection{Alternating Feedback}

\method{} improves A3 by 0.599~dB overall, with a 95\% image-pair-cluster
bootstrap interval of 0.551--0.648~dB. The mean gains are 0.680~dB under
stationary interference, 0.594~dB under alternating interference, and
0.523~dB under burst interference; their exploratory cluster intervals are
0.629--0.729, 0.540--0.649, and 0.469--0.578~dB, respectively. \method{}
exceeds A3 in 79.1\% of the 10,800 matched trials. Relative to A0, it gains
1.478~dB overall (95\% cluster interval: 1.392--1.563~dB).

The improvement is strongest at moderate and weak interference. After
pooling profiles, \method{} gains 0.405~dB at 0~dB SINR, 0.952~dB at 4~dB,
1.122~dB at 7~dB, and 0.882~dB at 10~dB. Table~\ref{tab:sinr_gain} also
reveals an important exception: under burst interference, updating the point
estimate is worse than retaining the A3 initialization at the two strongest
interference levels.

\begin{table}[t]
\caption{Paired PSNR Difference: \method{} $-$ A3 (dB)}
\label{tab:sinr_gain}
\centering
\setlength{\tabcolsep}{5.0pt}
\begin{tabular}{rrrrr}
\toprule
SINR & Stationary & Alternating & Burst & All \\
\midrule
$-7$ & $+0.323$ & $+0.206$ & $\mathbf{-0.174}$ & $+0.118$ \\
$-4$ & $+0.452$ & $+0.001$ & $\mathbf{-0.109}$ & $+0.115$ \\
$0$  & $+0.315$ & $+0.454$ & $+0.446$ & $+0.405$ \\
$4$  & $+1.084$ & $+0.925$ & $+0.846$ & $+0.952$ \\
$7$  & $+1.050$ & $+1.109$ & $+1.208$ & $+1.122$ \\
$10$ & $+0.854$ & $+0.870$ & $+0.923$ & $+0.882$ \\
\bottomrule
\end{tabular}
\end{table}

At burst SINR values of $-7$ and $-4$~dB, the losses are 0.174~dB
(95\% cluster interval: 0.126--0.224~dB) and 0.109~dB
(0.025--0.197~dB), respectively. Thus, diffusion-in-the-loop updates should
not always replace the initial local estimate.

\subsection{Reconstruction Versus Power Recovery}

The reconstruction gain is not explained by more accurate physical power
estimates. The A3 initializer has an aggregate block-power RMSE of 0.953, yet
only 28.3\% of \method{} trials finish with a smaller per-trial power RMSE.
At SINR values of 0, 4, and 10~dB, the median final power estimates are
0.234, 0.024, and 0.004, whereas the corresponding nominal mean powers are
0.990, 0.388, and 0.090. Fourteen trials contain at least one estimated block
power above 100; the maximum is 93,186.656.

Equation~\eqref{eq:update} should therefore not be described as recovering
the true nuisance parameter. Instead, the updates adapt the interference
scale and likelihood guidance along the reconstruction trajectory. Scale
ambiguity among $\vect{x}$, $\vect{z}$, and $a_b$, together with low-energy
interference estimates, may contribute to the discrepancy, but the historical
logs do not identify a single cause. The defensible conclusion is narrower:
reconstruction-optimal self-calibration need not coincide with physical
interference-power estimation.

\subsection{Runtime and Weak Interference}

The measured receiver times are 0.323~s per image for A3 and 0.347~s for
\method{} on an RTX~4090, a descriptive difference of 7.2\%. Because these
values come from separate batch-four runs, they are not a controlled overhead
measurement. The update itself adds no neural-network call; an interleaved
same-device benchmark is required for a final runtime claim.

At 10~dB SINR, \method{} is 0.882~dB better than A3 but remains 1.070~dB
below direct decoding (95\% cluster interval: $-1.270$ to $-0.867$~dB).
When interference is already weak, running a 40-step cancellation process can
therefore harm reconstruction. A deployable receiver should include a
validation-selected bypass based on receiver-observable reliability, not on
the unavailable true SINR.

\balance
\section{Discussion and Limitations}

The experiments separate three effects that are often conflated. First, true
global SINR is not needed to match the mean performance of the original global
receiver. Second, local power structure matters when interference is
nonstationary. Third, feeding evolving diffusion estimates back into
calibration improves reconstruction beyond one-shot local initialization.
The third effect survives all three profiles on average, but not every
operating point.

The negative results are part of the contribution. Under severe burst
interference, early point estimates are not reliable enough to improve the
block initializer. Under weak interference, diffusion cancellation itself is
unnecessary. These regimes suggest a three-way receiver decision among direct
decoding, global calibration, and local iterative cancellation. Designing
that decision from receiver observables is preferable to using an oracle SINR
threshold.

Several limitations bound the present claims. The channel is AWGN; the
receiver knows the noise level and block partition; only one desired-image
distribution, one interference-image distribution, one block length, and
fixed matched priors are evaluated; and PSNR is the only reconstruction
endpoint currently logged. LPIPS should be added to the independent
confirmation because pixel fidelity alone does not establish perceptual
improvement. Rayleigh fading, prior mismatch, randomized block boundaries,
and other block lengths remain untested. Finally, the reported image pairs
influenced method development. They are suitable for mechanism discovery but
not for the final confirmatory claim.

\section{Conclusion}

We removed true receiver-side SINR from diffusion interference cancellation
and studied the harder case in which interference power changes within a
frame. A blind global estimate preserved the mean performance of the
true-global-SINR receiver, while blockwise initialization captured gains that
a global parameter missed under alternating and burst interference. Closing
the loop between diffusion reconstruction and local amplitude calibration
added 0.599~dB over the one-shot blockwise receiver in the exploratory
evaluation. The same experiments exposed two operating boundaries: point
updates can hurt under severe burst interference, and diffusion cancellation
can hurt when interference is weak. Moreover, the reconstruction gain does
not imply accurate recovery of physical interference power. These findings
support self-calibrated, reliability-gated interference cancellation rather
than an unconditional receiver driven by oracle SINR.

\begingroup
\sloppy
\bibliographystyle{IEEEtran}
\bibliography{references}
\endgroup

\end{document}
